\author{OTS}
\documentclass[12pt]{mybib}
\setheaderlogo{logo.png}

% Absatzformatierung
\setlength{\parindent}{0pt}
\setlength{\parskip}{0.6em}
\renewcommand{\arraystretch}{1.8}

\title{2. Petrus}
\author{Lothar Schmid}
\date{}

\begin{document}
\mytitlepage
\tableofcontents

\section{2. Petrus 3}

\subsection{2. Petr 3,1-3}
 \begin{auslegung}{2. Petrus 3} 
        \hh{1}Geliebte, dies ist nun schon der zweite Brief, den ich euch schreibe, um durch Erinnerung eure lautere Gesinnung aufzuwecken,
        \switchcolumn
        \linkbib{IIPetr}{3:17}\\
        \linkbib{Rom}{11:28}\\
        \linkbib{Kol}{3:12}\\
        \linkbib{Kol}{4:14}
        \switchcolumn*
        \hh{2}damit ihr an die Worte gedenkt, die von den heiligen Propheten vorausgesagt worden sind und dessen, was euch der Herr und Retter durch uns, die Apostel, aufgetragen hat.
        \switchcolumn
        \linkbib{Lk}{24:6}\\
        \linkbib{Ps}{143:5}\\
        \linkbib{IIPetr}{1:19}\\
        \linkbib{IKor}{22:23}\\
        \linkbib{IThes}{2:13}\\
        \linkbib{Jud}{1:17}
        \switchcolumn*
        Hier die Erklärung
    \end{auslegung}




\end{document}